% ============================================================================ 6 · Producto
% fuente: docs/pitch/presentacion-contenido.md bloque 4 (principio, circuito, agentes, destinatarios, mensaje de ejemplo),
% feat/demo-gru:docs/memoria/f9-demo-gru-final.md (replay por punto de la perilla, 465 textos, 196 resúmenes,
% 0 plantillas), feat/demo-gru:PRODUCT.md (verificador, caché). Intentos rechazados: 103, el valor del pitch, por decisión
% del equipo (02-10); la ficha f9-demo-gru-final.md registra 40 (6 + 9 + 11 + 14).
\section{Producto: de la alerta a la acción}
\label{sec:producto}

Una alerta sola no evita ninguna falla: la evita que alguien haga algo a tiempo. El producto convierte cada alerta en una
acción ---un consejo de manejo al conductor o un contacto del concesionario--- redactada por agentes de IA generativa y
controlada por código. Lo ordena un principio: \emph{el modelo decide, el agente comunica y el código cuida lo que se dice}.
El modelo de lenguaje nunca produce un puntaje, un número ni un diagnóstico: solo redacta a partir de los hechos que el
sistema le entrega.

\subsection{El circuito}

\begin{figure}[!htbp]
\centering
\begin{tikzpicture}[
  node distance=4mm and 3.2mm,
  box/.style={draw,rounded corners=1.5pt,align=center,font=\footnotesize,minimum height=11mm,text width=21mm,inner sep=2pt},
  arr/.style={-{Stealth[length=2mm]},thick}]
\node[box] (t) {Telemetría\\viajes y señales};
\node[box,right=of t] (s) {Puntaje GRU\\cada 500 km};
\node[box,right=of s] (a) {Alerta\\2 cortes sobre $\tau_\alpha$};
\node[box,right=of a] (w) {Explicación\\hábitos vs.\ sanos};
\node[box,right=of w,dashed] (p) {Política\\determinista};
\node[box,below=6mm of p] (g) {Agentes LLM\\+ verificador};
\node[box,left=of g,text width=33mm] (u) {Conductor · concesionario\\· posventa};
\node[box,left=of u,text width=30mm] (r) {Retorno del taller\\monitoreo y reentrenamiento};
\draw[arr] (t) -- (s); \draw[arr] (s) -- (a); \draw[arr] (a) -- (w); \draw[arr] (w) -- (p);
\draw[arr] (p) -- (g); \draw[arr] (g) -- (u); \draw[arr] (u) -- (r); \draw[arr] (r) -| (t);
\end{tikzpicture}
\caption{Circuito de posventa. El modelo decide si hay alerta, la política (recuadro punteado) decide la acción y los
agentes solo redactan.}
\label{fig:circuito}
\end{figure}

El circuito de la Figura~\ref{fig:circuito} tiene siete pasos:
\begin{enumerate}
  \item \textbf{Telemetría diaria} del vehículo conectado: los resúmenes de viaje y los mensajes del filtro que Ford ya recibe.
  \item \textbf{Puntaje} de la GRU cada 500\,km, sobre los últimos 1.000\,km (\S\ref{sec:gru}).
  \item \textbf{Alerta confirmada} cuando el puntaje supera dos veces seguidas el umbral $\tau_\alpha$~\eqref{eq:alerta},
    con el presupuesto $\alpha$ que eligió Ford.
  \item \textbf{El porqué}: los hábitos en los que el vehículo supera al percentil 75 de los vehículos sanos de su mercado,
    solo los modificables y coherentes con la física del DPF, hasta tres (\S\ref{sec:explicabilidad}).
  \item \textbf{Política determinista}, que define Ford: si hay un hábito para nombrar, se envía un consejo al conductor; si
    no lo hay, lo contacta el concesionario; si el riesgo sigue alto dos revisiones después (una revisión es un corte, cada
    500\,km), el caso se escala al concesionario.
  \item \textbf{Redacción}: los agentes escriben los textos y un verificador programado los controla.
  \item \textbf{Retorno del taller}: cada intervención registrada indica si hubo alerta antes y con cuánto margen, y alimenta
    el monitoreo y el reentrenamiento.
\end{enumerate}

\subsection{Agentes y verificador}

La Tabla~\ref{tab:agentes} resume las tres piezas. Los agentes trabajan solo sobre las alertas, no sobre la flota, y
ninguno puede cambiar la acción que fija la política.

\begin{table}[!htbp]
\centering
\caption{Las piezas que redactan y controlan los textos.}
\label{tab:agentes}
\small
\begin{tabularx}{\textwidth}{@{}>{\raggedright\arraybackslash}p{3.4cm}>{\raggedright\arraybackslash}X@{}}
\toprule
Pieza & Qué hace \\
\midrule
Agente de clasificación & Recorre los eventos de la semana con herramientas: consulta la acción que fija la política y el
  contexto de la flota, encarga la redacción y arma el resumen semanal para posventa. No puede elegir una acción distinta de
  la política. \\
Agente redactor & Escribe el mensaje al conductor, el resumen para el taller y la línea de la bandeja de posventa. Las
  recomendaciones y los chequeos se eligen de listas cerradas, y su texto lo pone el código. \\
Verificador (código) & Todo número tiene que estar en los hechos del caso. Rechaza el lenguaje causal (``porque'',
  ``debido a''), las certezas, las promesas de que el riesgo baja, la palabra ``probabilidad'' y las atribuciones (``el
  sistema lo marcó por\ldots'', ``se parece a autos que fallaron''). Al conductor no le llegan síntomas del filtro (carga,
  regeneraciones, aceite, consumo). Si rechaza un texto, el agente lo reintenta hasta dos veces con la lista de problemas;
  si no lo logra, se usa una plantilla aprobada. \\
\bottomrule
\end{tabularx}
\end{table}

Cada pedido al modelo de lenguaje se guarda con una clave derivada de lo que lo define, de modo que un texto aprobado no
cambia entre ejecuciones y el prototipo se reproduce sin conexión.

\subsection{Tres destinatarios, tres mensajes}

\begin{itemize}
  \item \textbf{El conductor}, en la aplicación de su Ford, recibe un hábito concreto, algo que puede hacer el fin de semana,
    sin códigos de error, sin porcentajes y sin hablar de fallas. Por ejemplo: \emph{``Comparado con autos sanos de tu zona,
    tu camioneta hace muchos viajes cortos. Un tramo de ruta de 20 minutos por semana ayuda a que el filtro se limpie
    solo.''}
  \item \textbf{El concesionario} recibe la ficha del vehículo, con las señales del filtro comparadas con las de los
    vehículos sanos de su mercado y los chequeos sugeridos.
  \item \textbf{Posventa de Ford} recibe cada lunes una bandeja con las alertas de la semana priorizadas y ya redactadas,
    para aprobar o editar.
\end{itemize}

\subsection{Prototipo}
\label{sec:prototipo}

Se construyó una aplicación que reproduce semana a semana los 426 vehículos de desarrollo con el modelo finalista: en cada
semana el sistema solo conoce lo que se sabía ese día, y ningún vehículo del conjunto de prueba aparece en ella. Tiene tres
vistas ---la bandeja de posventa, la ficha del vehículo y el seguimiento de lo que pasó después--- y un control del
presupuesto $\alpha$ (Tabla~\ref{tab:replay}). El acceso está en la \S\ref{sec:complementaria}.

\begin{table}[!htbp]
\centering
\caption{La temporada reproducida en el prototipo, según el presupuesto de falsas alarmas \tagdev{} (426 vehículos: 135 con
evento y 291 sanos).}
\label{tab:replay}
\small
\begin{tabular}{@{}lrrrr@{}}
\toprule
 & $\alpha = 5\,\%$ & $10\,\%$ & $15\,\%$ & $20\,\%$ \\
\midrule
Vehículos con evento avisados antes (de 135) & 39 & 56 & 69 & 82 \\
Vehículos sanos con un aviso de más (de 291) & 14 & 29 & 43 & 58 \\
Alertas nuevas / escalamientos & 53 / 7 & 85 / 10 & 112 / 22 & 140 / 36 \\
Alertas con un hábito para nombrar & 48 & 73 & 101 & 125 \\
Anticipación mediana, semanas (km) & 15 (7.000) & 16 (7.200) & 16 (7.500) & 21 (8.100) \\
\bottomrule
\end{tabular}
\end{table}

Al $5\,\%$, 39 de los 135 vehículos con evento reciben un aviso con una mediana de 15 semanas de margen, y 14 de los 291
sanos reciben uno de más; 48 de las 53 alertas son un consejo al conductor. La muestra está enriquecida en eventos (135 de
426), de modo que los conteos no se trasladan directamente a una flota real (\S\ref{sec:costo}). En los cuatro puntos del
control, los agentes redactaron 465 textos y 196 resúmenes semanales, todos aprobados por el verificador, que rechazó 103
intentos en el camino; no hizo falta ninguna plantilla.
